% Source: https://artofproblemsolving.com/community/c3h1642667_inscribed_circle
% Stable identity: aops-topic-1642667-q01
Let $O$ be the centre of the regular hexagon. The vertices $N$ and $S$ are opposite, and the line $NS$ is perpendicular to the short diagonal $CA$ (both $C,A$ are symmetric about $NS$).

\textbf{$N'$ is the centre $O$.} Triangle $OCN$ is equilateral with side $1$, and $CA$ crosses the segment $ON$ at its midpoint (the distance from $O$ to $CA$ is $1\cdot\cos60^\circ=\dfrac{1}{2}$, and $ON=1$). So the reflection of $N$ in $CA$ is $N'=O$.

\textbf{The circle.} A circle tangent to $ST$ and $SM$ has its centre $P$ on the bisector of $\angle TSM=120^\circ$, which is the line $SO$. If $SP=d$, its radius is the distance from $P$ to $SM$:
\[
\rho=d\sin60^\circ=\dfrac{\sqrt3}{2}d .
\]

\begin{asy}
usepackage("amsmath");
size(6.5cm);
defaultpen(fontsize(10pt));
real s3=sqrt(3);
pair Ms=dir(330),Cs=dir(30),Ns=dir(90),As=dir(150),Ts=dir(210),Ss=dir(270),O=(0,0);
real d=2*(2-s3), r=2*s3-3; pair P=Ss+(0,d);
draw(Cs--Ns--As,gray+dashed);draw(Cs--As,gray+dashed);
draw(Ms--Cs--O--As--Ts--Ss--cycle,linewidth(0.9));
draw(circle(P,r),red+linewidth(0.9));
pair F1=Ss+(d*cos(pi/3))*unit(Ms-Ss),F2=Ss+(d*cos(pi/3))*unit(Ts-Ss);
draw(P--F1^^P--F2^^P--O,blue+dashed);dot(P);dot(O);dot(Ns);
label("$M$",Ms,SE);label("$C$",Cs,NE);label("$N$",Ns,N);label("$A$",As,NW);label("$T$",Ts,SW);label("$S$",Ss,S);
label("$N'$",O+(0,0.06),N);label("$P$",P,dir(20));
label("\mbox{}",(min(currentpicture,user=true)+max(currentpicture,user=true))/2,black); // keeps the text after the figure black when LaTeX typesets the labels inline
\end{asy}

It passes through $O$, and $SO=1$. If $P$ lies between $S$ and $O$, then $PO=1-d=\rho$:
\[
1-d=\dfrac{\sqrt3}{2}d\quad\Longrightarrow\quad d=\dfrac{2}{2+\sqrt3}=2(2-\sqrt3),\qquad \rho=\sqrt3(2-\sqrt3)=2\sqrt3-3\approx0.464 .
\]
(If $P$ were beyond $O$, then $d-1=\dfrac{\sqrt3}{2}d$ gives $d=2(2+\sqrt3)>2$, far outside the hexagon.) This circle lies in the concave hexagon: it touches $SM$ and $ST$ at points of these sides (at distance $d\cos60^\circ=2-\sqrt3<1$ from $S$), and at $N'=O$ its tangent line is perpendicular to $SO$, while the sides $N'C$ and $N'A$ point away from $S$ (each makes $120^\circ$ with $OS$), so it does not cross them.

The area is
\[
\pi\rho^2=\pi\left(2\sqrt3-3\right)^2=\boxed{(21-12\sqrt3)\pi}\approx0.676 .
\]
(For side length $s$ the area is $(21-12\sqrt3)\pi s^2$.)
